\problem{3.23}{Projection operators and their spectrum}
For a normalized ket $\ket\alpha$ define
$P=\ket\alpha\bra\alpha$. Then
\begin{equation}
P^2=\ket\alpha\braket\alpha\alpha\bra\alpha
=\ket\alpha\bra\alpha=\boxed P.
\end{equation}
Any eigenvalue $p$ obeys $p^2=p$, hence
$p=0$ or $1$. The vector $\ket\alpha$
is a $p=1$ eigenvector, and every vector
orthogonal to it is in the $p=0$
eigenspace. A projection onto a
higher-dimensional subspace has the same
eigenvalues, with the range and kernel
as the corresponding eigenspaces.

\problem{3.24}{Hermiticity of an operator's matrix}
Let $\{\ket n\}$ be an orthonormal basis and
$Q_{mn}=\bra m Q\ket n$.
If $Q=Q^\dagger$,
\begin{align}
Q_{mn}&=\braket m{Qn}
=\braket{Qm}n\\
&=\braket n{Qm}^*
=\boxed{Q_{nm}^*}.
\end{align}
Thus the matrix $Q$ satisfies
$\boxed{Q=Q^\dagger}$, i.e.\ its
conjugate transpose equals itself.
In particular all diagonal entries are real.

\problem{3.25}{Diagonalizing a two-level Hamiltonian}
In the orthonormal basis $\ket1,\ket2$,
the stated operator has the matrix
\begin{equation}
H=\varepsilon
\begin{pmatrix}1&1\\1&-1\end{pmatrix}.
\end{equation}
Its characteristic equation is
\begin{equation}
\det(H-EI)=E^2-2\varepsilon^2=0,
\qquad\boxed{E_\pm=\pm\sqrt2\,\varepsilon}.
\end{equation}
For $E_+$ the first row gives
$v_2=(\sqrt2-1)v_1$; for $E_-$,
$v_2=-(\sqrt2+1)v_1$.
A normalized eigenbasis is therefore
\begin{align}
\boxed{\ket{E_+}}&=
\frac{\ket1+(\sqrt2-1)\ket2}
{\sqrt{1+(\sqrt2-1)^2}},\\
\boxed{\ket{E_-}}&=
\frac{\ket1-(\sqrt2+1)\ket2}
{\sqrt{1+(\sqrt2+1)^2}}.
\end{align}
The vectors are orthogonal because
$1+(\sqrt2-1)(-\sqrt2-1)=0$.

\problem{3.26}{Complex kets, dual vectors and outer products}
With basis coordinates
$\ket\alpha=(\ii,-2,-\ii)^{\mathsf T}$,
$\ket\beta=(\ii,0,2)^{\mathsf T}$,
taking conjugate transposes gives
\begin{equation}
\boxed{\bra\alpha=(-\ii,-2,\ii)},\qquad
\boxed{\bra\beta=(-\ii,0,2)}.
\end{equation}
The two inner products are
\begin{equation}
\boxed{\braket\alpha\beta=
(-\ii)\ii+\ii(2)=1+2\ii},\qquad
\boxed{\braket\beta\alpha=1-2\ii}.
\end{equation}
The outer product $A=\ket\alpha\bra\beta$ is
\begin{equation}
\boxed{A=
\begin{pmatrix}
1&0&2\ii\\
2\ii&0&-4\\
-1&0&-2\ii
\end{pmatrix}}.
\end{equation}
For example $A_{13}=\ii\cdot2=2\ii$ and
$A_{31}=(-\ii)(-\ii)=-1$.
Since $A_{13}\ne A_{31}^*$,
$\boxed{A\ne A^\dagger}$.

\problem{3.27}{Spectral decomposition and functions of an operator}
Let $Q\ket{e_n}=q_n\ket{e_n}$, where
$\{\ket{e_n}\}$ is a complete orthonormal
basis. Completeness gives
$\sum_n\ket{e_n}\bra{e_n}=I$. On an arbitrary
vector $\ket\alpha=\sum_n\alpha_n\ket{e_n}$,
\begin{equation}
Q\ket\alpha=
\sum_nq_n\alpha_n\ket{e_n}=
\left(\sum_nq_n\ket{e_n}\bra{e_n}\right)\ket\alpha.
\end{equation}
Therefore
\begin{equation}
\boxed{Q=\sum_nq_n\ket{e_n}\bra{e_n}}.
\end{equation}
For a function $f$ defined on the spectrum,
$f(Q)=\sum_nf(q_n)\ket{e_n}\bra{e_n}$.
For the exponential, its power series gives
\begin{align}
\ee^Q&=\sum_{k=0}^\infty Q^k/k!\\
&=\sum_n\sum_{k=0}^\infty
q_n^k/k!\ket{e_n}\bra{e_n}\\
&=\boxed{\sum_n\ee^{q_n}\ket{e_n}\bra{e_n}}.
\end{align}
\qnote{For unbounded operators, convergence and
the domain of each function of $Q$ are
specified by the spectral theorem.}

\problem{3.28}{Functions of the differentiation operator}
For $D=\dd/\dd x$, apply
$\sin D=\sum_{j\ge0}(-1)^jD^{2j+1}/(2j+1)!$
to $x^5$. Higher than fifth derivatives vanish:
\begin{equation}
\boxed{(\sin D)x^5=
5x^4-\frac{60x^2}{6}+\frac{120}{120}
=5x^4-10x^2+1}.
\end{equation}
To evaluate $(1-D/2)^{-1}\cos x$, solve
$(1-D/2)f=\cos x$. Set
$f=a\cos x+b\sin x$.
Since $D\cos x=-\sin x$ and
$D\sin x=\cos x$,
\begin{equation}
(a-b/2)\cos x+(b+a/2)\sin x=\cos x.
\end{equation}
The equations $a-b/2=1$, $b+a/2=0$
give $a=4/5$ and $b=-2/5$, hence
\begin{equation}
\boxed{(1-D/2)^{-1}\cos x=
\frac45\cos x-\frac25\sin x}.
\end{equation}
\qnote{The inverse is interpreted on the
two-dimensional invariant span of
$\cos x$ and $\sin x$, where it is
unambiguously invertible.}

\problem{3.29}{A central-commutator Baker–Campbell–Hausdorff formula}
Assume $C=[A,B]$ and $[A,C]=[B,C]=0$.
Induction from
$[A^{n+1},B]=A[A^n,B]+[A,B]A^n$
gives
\begin{equation}
\boxed{[A^n,B]=nA^{n-1}C}.
\end{equation}
Expanding $\ee^{\lambda A}$ into a series,
\begin{align}
[\ee^{\lambda A},B]
&=\sum_{n=1}^\infty\frac{\lambda^n}{n!}
[A^n,B]\\
&=\lambda\ee^{\lambda A}C.
\end{align}
To establish the BCH identity, define
$F(t)=\ee^{t(A+B)}$ and
$G(t)=\ee^{tA}\ee^{tB}\ee^{-t^2C/2}$.
Using $[A,\ee^{tB}]=tC\ee^{tB}$,
or equivalently
$\ee^{tA}B=(B+tC)\ee^{tA}$,
differentiate $G$:
\begin{equation}
G'=(A+B+tC-tC)G=(A+B)G.
\end{equation}
Since $G(0)=I=F(0)$ and
$F'=(A+B)F$, uniqueness gives
$F(t)=G(t)$ and at $t=1$,
\begin{equation}
\boxed{\ee^{A+B}=\ee^A\ee^B\ee^{-C/2}}.
\end{equation}
\qnote{For unbounded operators, this algebra
is formal unless the exponential operators
and their domains satisfy the necessary
analytic conditions. The Weyl relations
provide a rigorous form for $x,p$.}

\problem{3.30}{From position-space to energy-space amplitudes}
Let $\ket n$ be a discrete orthonormal
energy eigenbasis, with position representation
$\psi_n(x)=\braket x n$. For a state $\ket{S(t)}$,
insert the position resolution of the identity:
\begin{align}
c_n(t)&=\braket n{S(t)}
=\bra n\left(\int_{\mathbb R}
\ket x\bra x\dd x\right)\ket{S(t)}\\
&=\int_{\mathbb R}
\braket n x\,\braket x{S(t)}\,\dd x\\
&=\boxed{\int_{\mathbb R}
\psi_n^*(x)\Psi(x,t)\dd x}.
\end{align}
For a time-independent Hamiltonian,
\begin{equation}
\boxed{c_n(t)=c_n(0)\ee^{-\ii E_nt/\hbar}},
\quad\boxed{\Psi(x,t)=\sum_nc_n(t)\psi_n(x)}.
\end{equation}
This is the energy-basis transformation
of the same Hilbert-space vector.
